%BeginMC
\question
Lower elevations on a very high mountain near the equator experience more precipitation and warmer temperatures than high elevations. If you climbed from the base of the mountain to the top, you would notice a change in the types of ecosystems present on the mountain side. This would be similar to the changes (Note: The equator is 0°, the poles are at 90°). 

\begin{choices}
\choice traveling across South American from 55°S latitude to 5°S latitude.%EndChoice
\correctchoice traveling across North America from 15°N latitude to 55°N latitude.%EndChoice
\choice traveling across South America and North American from 55°S latitude to 55°N latitude.%EndChoice
\choice across Europe and Asia from 15°E longitude to 15°W longitude.%EndChoice
\choice across Europe and Asia from 15°W longitude to 15°E longitude.%EndChoice
\end{choices}
%EndMC

%BeginMC
\question
Considering your knowledge of the dead zone and the causes of microscopic algae blooms, the Gulf of Mexico ecosystem is \emph{most likely} to be limited by which of the following nutrients?

\begin{choices}
	\choice carbon%EndChoice
	\choice \ce{CO2}%EndChoice
	\choice hydrogen%EndChoice
	\choice oxygen%EndChoice
	\correctchoice nitrogen%EndChoice
\end{choices}
%EndMC


%BeginMC
\question
If humans stopped fishing along the coast of Alaska, what will probably happen to the kelp forest and why?

\begin{choices}
	\correctchoice The kelp abundance will increase because the sea urchin abundance will probably decrease.%EndChoice
	\choice The kelp abundance will remain about the same. The kelp are not part of this community, so the kelp would not be affected.%EndChoice
	\choice The kelp abundance will decrease because more fish will be present to eat the kelp.%EndChoice
	\choice The kelp abundance will decrease. Sea otters eat urchins and urchins eat kelp. Because there were fewer otters to eat urchins, the number of urchins increased. The urchins therefore ate more kelp so the kelp abundance decreased.%EndChoice
	\choice The kelp abundance will increase because of the symbiotic relationship between kelp and fishes. Fishes are necessary for kelp to reproduce.%EndChoice
\end{choices}
%EndMC

%BeginMC
\question
Which of these species is most likely a keystone species in Alaska?

\begin{choices}
	\choice Humans.  %EndChoice
	\choice Kelp.  %EndChoice
	\choice Killer whales. %EndChoice
	\correctchoice Sea otters. %EndChoice
	\choice Sea urchins.  %EndChoice
\end{choices}
%EndMC

%BeginMC
\question
The Mediterranean coast and coastal California have similar types of ecosystems. This is because

\begin{choices}
\correctchoice they have similar climates.%EndChoice
\choice they have similar types of organisms.%EndChoice
\choice they are both wine-growing regions.%EndChoice
\choice they have similar types of communities.%EndChoice
\choice they are at similar latitudes.%EndChoice
\end{choices}
%EndMC

%BeginMC
\question
The land on the south side of the Himalayan mountain range (below) receives relatively high levels of rainfall and has fertile farmlands.  The land north of the mountain range is hot and dry and can't be farmed.  Which of the following is true.

\begin{multicols}{2}
\begin{choices}
\choice The south side of the Himalayan mountain range is in a rain shadow.%EndChoice
\choice The desert in southern China warms the north side of the mountain range, causing it to dry out.%EndChoice
\choice The cool, moist Indian Ocean air heats up as it rises, releasing its precipitation as it passes the tops of the mountains, and this warm, now dry air cools as it descends on the leeward side of the range.%EndChoice
\correctchoice Warm, moist air from the Indian Ocean blows north, rises and cools, releasing precipitation as it moves up the south side of the range, and this cool, now dry air mass heats up as it descends on the north side of the range.%EndChoice
\choice The warm, moist air is too dense to rise over the height of Mount Everest and the other mountains in this range so the moisture remains on the south side of the range.%EndChoice
\end{choices}
\columnbreak
	\includegraphics[width=0.45\textwidth]{exam4_himalayas}
\end{multicols}
%EndMC

%BeginMC
\question
The Ozark region of Missouri lies within the temperate deciduous forest ecosystem.  If, within the next 25 years, global climate change raises the mean annual temperature from 14°C to 18°C and lowers the mean annual precipitation from 160 cm to 130 cm, which of the following is most likely?

\begin{multicols}{2}
	\begin{choices}
		\choice The ecosystem will remain temperature deciduous ecosystem. %EndChoice
		\choice The ecosystem will become more like boreal forest. %EndChoice
		\choice The ecosystem will become more like temperate rain forest. %EndChoice
		\choice The ecosystem will become more like temperate grassland. %EndChoice
		\correctchoice The ecosystem will become more like sclerophyllous or thorn woodland. %EndChoice
	\end{choices}
\columnbreak
	\includegraphics[width=0.45\textwidth]{exam4_climograph2}
\end{multicols}

%EndMC

%BeginMC
\question
On average, the amount of secondary consumer biomass is about \underline{\hspace{2cm}} less than the amount of primary producer biomass. 
\begin{choices}
\choice 10$\times$ %dontshuffle %EndChoice
\choice 20$\times$ %dontshuffle %EndChoice
\choice 50$\times$ %dontshuffle %EndChoice
\choice 80$\times$ %dontshuffle %EndChoice
\correctchoice 100$\times$%dontshuffle %EndChoice 
\end{choices}
%EndMC

%BeginMC
\question
A fish ate shrimp. The shrimp ate algae (which are photosynthetic).  The algae is a \rule{1in}{0.4pt}. The fish is a \rule{1in}{0.4pt}.

\begin{choices}	
	\choice primary consumer; secondary consumer. %EndChoice
	\choice secondary consumer; primary consumer. %EndChoice
	\choice secondary consumer; primary producer. %EndChoice
	\choice primary consumer; primary producer. %EndChoice
	\correctchoice primary producer; secondary consumer. %EndChoice
	\choice primary producer; primary consumer. %EndChoice
\end{choices}
%EndMC




%BeginMC
\question
Consider the food chain grass $\to$ grasshopper $\to$ mouse $\to$ snake $\to$ hawk. How much of the chemical energy fixed by photosynthesis of the grass (100\%) is available to the hawk?

\begin{choices}
\correctchoice 0.01\%%EndChoice
\choice 0.1\% %dontshuffle %EndChoice
\choice 1\% %dontshuffle %EndChoice
\choice 10\% %dontshuffle %EndChoice
\choice 60\% %dontshuffle %EndChoice
\end{choices}
%EndMC




%BeginMC
\question
Plants produce water vapor during photosynthesis, which goes from the plant to the atmosphere. The water vapor in the atmosphere eventually falls back to the earth as rain. This process would be part of

\begin{choices}
\correctchoice the water cycle.%EndChoice
\choice energy flow.%EndChoice
\choice plant metabolism.%EndChoice
\choice the entire ecosystem.%EndChoice
\choice the energy cycle.%EndChoice
\choice the nutrient cycle.%EndChoice
\end{choices}
%EndMC

%BeginMC
\question
The energy stored by producers can be passed through an ecosystem along a(n) \rule{1in}{0.4pt}, a series of interactions in which organisms transfer energy by eating and being eaten.

\begin{choices}
	\correctchoice food web %EndChoice
	\choice ecosystem %EndChoice
	\choice trophic cascade %EndChoice
	\choice energy cycle %EndChoice
	\choice nutrient cycle %EndChoice
	\choice predator-prey chain  %EndChoice
\end{choices}
%EndMC

%BeginMC
\question
In ecosystems, why is ``cycling'' used to describe the transfer of chemical elements and nutrients, while ``flow'' is used for the transfer of energy?

\begin{choices}
\correctchoice Chemical elements are recycled within ecosystems, but energy flows through and out of ecosystems.%EndChoice
\choice Both chemical elements and energy are recycled but can flow to other ecosystems.%EndChoice
\choice Chemical elements are cycled into ecosystems from other ecosystems, but energy constantly flows within the ecosystem.%EndChoice
\choice Both chemical elements and energy flow within an ecosystem but can cycle to other ecosystems.%EndChoice
\choice Energy flows from one ecosystem to other ecosystems, but chemical elements constantly cycle within an ecosystem.%EndChoice
\end{choices}
%EndMC

%BeginMC
\question
The dead zone in the Gulf of Mexico and other oceans around the globe occurs when

\begin{choices}
\choice an excess of nutrients causes an algal bloom. The algae consume all of the oxygen in the water for photosynthesis, so other organisms can no longer live there.%EndChoice
\choice eutrophication causes an absence of nutrients in the zone.%EndChoice
\choice commercial fishers over-harvest shrimp, crab, fishes and other sea animals.%EndChoice
\correctchoice an excess of nutrients causes an algal bloom, followed by death of the algae. Decomposition of the algae uses up the oxygen in the water, so other organisms can no longer live there.%EndChoice
\choice global climate change causes the water to become too warm to sustain life.%EndChoice
\end{choices}
%EndMC

%BeginMC
\question
A tertiary consumer has high abundance. In a trophic cascade, which of the following is probably \emph{not} true.
\begin{choices}
	\correctchoice The primary producer will have high abundance. %EndChoice
	\choice The secondary consumer will have low abundance. %EndChoice
	\choice The primary consumer will have high abundance. %EndChoice
	\choice The decomposer level will not be affected. %EndChoice%
\end{choices}


%BeginMC
\question
Based on the hypothetical food web below, which letter most likely represents a decomposer. Arrows point in the direction of energy flow.

\begin{multicols}{3} 
	\begin{choices}
		\correctchoice A %dontshuffle %EndChoice
		\choice B %dontshuffle %EndChoice
		\choice C %dontshuffle %EndChoice
		\choice D %dontshuffle %EndChoice
		\correctchoice E %dontshuffle %EndChoice
	\end{choices}
	\columnbreak
	\includegraphics{exam4_foodweb}
	\columnbreak
\end{multicols}
%EndMC

%BeginMC
\question
In 2011, the Mississippi River experienced a record flood event.  Which of the following statements is the best prediction of how this record flood event will affected the Gulf of Mexico dead zone that summer?

\begin{choices}
	\correctchoice The dead zone will be larger because the volume of water will carry even more nutrients into the Gulf of Mexico.%EndChoice
	\choice The dead zone will be smaller because the volume of water will bring more oxygen to the Gulf of Mexico, which will offset the hypoxia.%EndChoice
	\choice The dead zone will be smaller because the volume of water will dilute the nutrients, reducing eutrophication.%EndChoice
	\choice The dead zone will be similar to previous years because most of the flood water will soak into the ground and not enter the Gulf of Mexico.%EndChoice
	\choice The dead zone will be larger because the volume of water will cause more fishes to avoid the freshwater input.%EndChoice
\end{choices}
%EndMC

